\scanpage{74}
\mspage{11}
% archive foliation: 76
d) Les séries convergentes
\[
f=\sum_{i\geq0}a_i z^i\qquad(a_i\in\mathbb C)
\]
de rayon de convergence $\geq1$ (donc
définissant $\mathbb D^\circ\longrightarrow\mathbb C$)
\inserted{avec $|f(z)|\leq1$ si $z\in\mathbb D$}, telles que la $f$
se prolonge par continuité en une
fonction $\mathbb D\longrightarrow\mathbb D$, celle-ci ayant les
propriétés 1) 2) 3) de B) ci-dessus.

\textcircled{6}
\margin{\uline{13.1.}}
Soit $(X,C)$ un objet du groupoïde
de Teichmüller analytique complexe
à bord, avec la condition différentiable.
Pour fixer les idées, supposons $X$ connexe,
de type $g,\nu$. Considérons la découpe \inserted{$X'$}
de $X$ suivant $C$, bouchons ses trous
\inserted{(canoniquement en utilisant la structure torique de $\delta X'$)*},
\margin{* via la métrique riemannienne can. de $X'$}
d'où \unsure{surfaces} holomorphes compactes \inserted{sans bord $X''$}, à disques
\inserted{ponctués} marqués, avec des \deleted{iso}\inserted{\unsure{antiiso|antiso}} de \deleted{ces}\inserted{certains de ces} disques
(qui se groupent \deleted{deux à deux} par paires, à l'exception
des $\nu$ disques correspondant à $\delta X$).
On peut considérer $X''$ comme la
normalisée d'une courbe \deleted{$\unsure{R}$}\inserted{$\mathbb C$-alg $X'''$} stable
(au sens de Mumford-Deligne), en
